\section{Introduction to Probability}

\emph{Probability} is the branch of mathematics that studies random phenomena—that is, phenomena whose outcomes cannot be determined in advance with certainty. Unlike deterministic phenomena (such as an object falling under gravity), random phenomena exhibit inherent variability that can only be described using probabilistic models.

\subsection{Why study probability?}

Probability is fundamental across numerous areas of knowledge:

\begin{itemize}
	\item \textbf{Natural sciences:} Modeling of quantum phenomena, genetics, evolution, and statistical thermodynamics.
	\item \textbf{Engineering:} System reliability, queueing theory, and signal processing.
	\item \textbf{Economics and finance:} Option pricing, risk management, and market models.
	\item \textbf{Computer science:} Randomized algorithms, machine learning, and cryptography.
	\item \textbf{Medicine:} Clinical trials, epidemiology, and probabilistic diagnosis.
	\item \textbf{Daily life:} Decision-making under uncertainty, games of chance, and weather forecasting.
\end{itemize}

\subsection{Relationship with statistics}

Probability and statistics are closely intertwined:

\begin{itemize}
	\item \textbf{Probability} provides the theoretical framework for modeling random phenomena and calculating the likelihood of events.
	\item \textbf{Inferential statistics} uses observed data to make inferences about populations, drawing on probabilistic models as its foundation.
\end{itemize}

In other words, probability tells us how likely it is to observe certain data given a model, whereas statistics helps us infer the model from the observed data.

\subsection{Fundamental concepts}

In this chapter, we will develop the following concepts:

\begin{enumerate}
	\item \textbf{Random experiments and sample space:} We will define what a random experiment is and how to describe all its possible outcomes.
	\item \textbf{Events and set operations:} We will study how to combine events using union, intersection, and complement.
	\item \textbf{Axioms of probability:} We will present Kolmogorov's modern axiomatic approach.
	\item \textbf{Conditional probability and independence:} We will analyze how the occurrence of one event affects the probability of another.
	\item \textbf{Bayes' Theorem:} A fundamental tool for updating probabilities with new information.
	\item \textbf{Random variables:} Functions that assign numerical values to the outcomes of random experiments.
	\item \textbf{Special distributions:} Commonly used probabilistic models (binomial, normal, Poisson).
\end{enumerate}

\subsection{A brief history}

The mathematical study of probability began in the 17th century with correspondence between Blaise Pascal and Pierre de Fermat regarding problems in games of chance. Later, Jacob Bernoulli proved the Law of Large Numbers, and in the 20th century, Andrey Kolmogorov established the modern axioms of probability in 1933, providing a rigorous mathematical foundation.

\begin{remark}
	Throughout this chapter, we will use concrete examples (tossing coins, rolling dice, drawing cards) to illustrate concepts. However, these exact same principles apply to much more complex problems across science, engineering, and decision-making.
\end{remark}
